\documentclass[11pt]{article}
\usepackage[margin=0.82in]{geometry}
\usepackage{fontspec}
\setmainfont{STIXGeneral}
\usepackage{microtype}
\usepackage{xcolor}
\usepackage[colorlinks=true,urlcolor=blue!45!black]{hyperref}
\hypersetup{pdftitle={Cover Letter: Testing Singular Boundary-State Sufficiency in Gowdy Cosmology},pdfauthor={Ronald Bibb}}
\pagestyle{empty}
\setlength{\parskip}{5pt}
\setlength{\parindent}{0pt}

\begin{document}

\begin{flushright}
Ronald Bibb\\
Independent Researcher\\
Lilburn, GA 30047, USA\\
\href{mailto:ronbibb@gmail.com}{ronbibb@gmail.com}\\
ORCID: \href{https://orcid.org/0009-0004-1153-2464}{0009-0004-1153-2464}\\[7pt]
10 August 2026
\end{flushright}

Editors\\
\textit{Classical and Quantum Gravity}

\vspace{7pt}

Dear Editors,

Please consider my manuscript, ``Testing Singular Boundary-State Sufficiency
in Gowdy Cosmology: Smooth Closure and Spike-Localized Failure,'' for publication as a regular
research article in \textit{Classical and Quantum Gravity}.

The manuscript reverses the usual inferential direction of a singular
initial-value calculation. Starting from ordinary finite-time Cauchy data, it
extracts a proposed asymptotic state, discards the bulk history, independently
reconstructs a second branch, and tests whether the proposed state preserves
declared geometric observables.

The principal result is a nonlinear negative result. In generic unpolarized Gowdy
$T^3$ data, ordinary four-function velocity-dominated data produce a localized
reconstruction distance near $0.82$ at all three resolutions around two
resolved spike candidates. Exact true-spike controls explain the failure:
ordinary data lose the spike, whereas an enriched centered-spike state restores
the exact family below $2.4\times10^{-4}$. The manuscript does not claim generic
almost-everywhere sufficiency.

Smooth polarized Einstein--scalar data and an exact Bessel benchmark supply
the controlled positive comparison. They converge at second order and show
that normalized shape and full marked geometry require different data. The
additive metric mode is analytically familiar and is not claimed as a new
degree of freedom. Its role is supporting evidence; the distinguishing result
is that the same diagnostic closes in the smooth sector and rejects an
insufficient state in the exceptional nonlinear sector.

The exceptional status of Gowdy spikes is established in the analytical
literature and is not presented as new. The contribution is instead a
quantitative, observable-relative reconstruction-loss diagnostic and an exact
control identification of the enriched data that repair that loss.

The extraction-depth logic is situated alongside finite-radius and
Cauchy-characteristic waveform extraction in numerical relativity. The Gowdy
asymptotic and spike literature engaged by the manuscript has appeared
substantially in
\textit{Classical and Quantum Gravity}, making the journal a natural venue.
The versioned source, machine-readable outputs, exact controls, independent
verification records, and stopped deep-run outcomes are publicly available in
the repositories identified in the manuscript.

The manuscript is original and is not under consideration elsewhere.

Thank you for considering this manuscript.

Sincerely,

\vspace{14pt}

Ronald Bibb\\
Independent Researcher\\
\href{mailto:ronbibb@gmail.com}{ronbibb@gmail.com}\\
ORCID: \href{https://orcid.org/0009-0004-1153-2464}{0009-0004-1153-2464}

\end{document}
